\section{Benchmarks and Experimental Design}
\label{sec:benchmarks}

\subsection{Benchmark Datasets}

We use the standard SR benchmarks from Liu et~al.\ \cite{liu2025} (Table~1)
and Udrescu \& Tegmark \cite{udrescu2020}.

\subsubsection{Nguyen Benchmarks (12 expressions)}

\begin{table}[h]
\centering
\caption{Nguyen symbolic regression benchmarks (from Liu2025, Table 1).}
\label{tab:nguyen}
\begin{tabular}{llcl}
\toprule
ID & Expression & Vars & Range \\
\midrule
Nguyen-1  & $x^3 + x^2 + x$ & 1 & $[-1, 1]$ \\
Nguyen-2  & $x^4 + x^3 + x^2 + x$ & 1 & $[-1, 1]$ \\
Nguyen-3  & $x^5 + x^4 + x^3 + x^2 + x$ & 1 & $[-1, 1]$ \\
Nguyen-4  & $x^6 + x^5 + x^4 + x^3 + x^2 + x$ & 1 & $[-1, 1]$ \\
Nguyen-5  & $\sin(x^2)\cos(x) - 1$ & 1 & $[-1, 1]$ \\
Nguyen-6  & $\sin(x) + \sin(x + x^2)$ & 1 & $[-1, 1]$ \\
Nguyen-7  & $\log(x+1) + \log(x^2+1)$ & 1 & $[0, 2]$ \\
Nguyen-8  & $\sqrt{x}$ & 1 & $[0, 4]$ \\
Nguyen-9  & $\sin(x) + \sin(y^2)$ & 2 & $[-1, 1]$ \\
Nguyen-10 & $2\sin(x)\cos(y)$ & 2 & $[-1, 1]$ \\
Nguyen-11 & $x^y$ & 2 & $[0, 1]$ \\
Nguyen-12 & $x^4 - x^3 + \frac{y^2}{2} - y$ & 2 & $[-1, 1]$ \\
\bottomrule
\end{tabular}
\end{table}

\textbf{Data configuration} (following Liu2025 Section 4.1): 20 training
points and 100 test points uniformly sampled from the specified range.
All experiments use seed 42 for reproducibility.

\subsubsection{Feynman Physics Benchmarks}

10 selected equations from the AI Feynman dataset \cite{udrescu2020},
as used in Liu2025 Table 2. These range from simple products
(e.g., $F = \mu N_s$) to relativistic expressions
(e.g., $m = m_0 / \sqrt{1 - v^2/c^2}$), with 1--3 variables each.

\subsection{Evaluation Metrics}

\begin{itemize}[nosep]
  \item \textbf{$R^2$ (Coefficient of Determination)}: Primary metric.
    $R^2 = 1 - \frac{\sum_i (y_i - \hat{y}_i)^2}{\sum_i (y_i - \bar{y})^2}$.
    Perfect fit: $R^2 = 1$.
  \item \textbf{NRMSE}: $\text{NRMSE} = \frac{\sqrt{\text{MSE}}}{\sigma_y}$.
    Scale-invariant error metric.
  \item \textbf{Reward} (Liu2025, Eq.~12):
    $R(\mathcal{G}) = \frac{1}{\text{NRMSE} + 1}$.
    Bounded in $(0, 1]$.
\end{itemize}

\subsection{Experimental Design: WITH vs WITHOUT Canonicalization}

The key experiment compares the \textbf{same search algorithm} operating
in two modes:
\begin{enumerate}[nosep]
  \item \textbf{Canonical mode}: Every generated/mutated string is canonicalized
    before evaluation. Duplicates detected via canonical string equality
    ($O(k!)$ deduplication).
  \item \textbf{Non-canonical mode} (\texttt{--no-canon}): Strings are evaluated
    as-is, without canonicalization. No deduplication.
\end{enumerate}

This paired design directly measures the benefit of the $O(k!)$ search space
reduction without confounding it with algorithm differences.

\subsection{Search Space Reduction Analysis}

The primary experiment generates $N$ random IsalSR strings, computes the
canonical string for each, and measures:
\begin{equation}
  \text{Reduction factor} = \frac{N_{\text{valid}}}{N_{\text{unique canonical}}}
\end{equation}

This is compared against the theoretical bound $k!$ where $k$ is the average
number of internal nodes. A reduction factor approaching $k!$ validates the
paper's central claim.
